% ============================================================================
%  Fourier Series and the Discrete Fourier Transform (DFT)
%  Computational Physics II - Unit 4: Fourier Transform I
%  Universidad de Pamplona - Faculty of Basic Sciences
%  Compile with:  pdflatex fourier_slides.tex   (twice)
%  Needs in the same folder (generated by running  python fourier_physics.py):
%    fig1_partial_sums.png, fig2_normal_modes.png, fig3_spectral_leakage.png,
%    fig4_decay_rates.png,  fig5_zero_padding.png
% ============================================================================
\documentclass[aspectratio=169,10pt]{beamer}

\usepackage[utf8]{inputenc}
\usepackage[T1]{fontenc}
\usepackage[english]{babel}
\usepackage{amsmath,amssymb,amsfonts}
\usepackage{graphicx}
\usepackage[most]{tcolorbox}
\usepackage{booktabs}

% ---------------------------------------------------------------------------
% Wine-red (vinotinto) color scheme -- edit the RGB values here
% ---------------------------------------------------------------------------
\definecolor{wine}{RGB}{114,28,48}
\definecolor{winelight}{RGB}{246,232,235}
\definecolor{winedark}{RGB}{78,16,32}
\definecolor{gold}{RGB}{176,125,35}
\definecolor{softgray}{RGB}{95,85,88}

\setbeamercolor{frametitle}{bg=wine,fg=white}
\setbeamercolor{title}{fg=wine}
\setbeamercolor{structure}{fg=wine}
\setbeamercolor{itemize item}{fg=wine}
\setbeamercolor{itemize subitem}{fg=gold}
\setbeamercolor{enumerate item}{fg=wine}
\setbeamercolor{normal text}{fg=black}
\setbeamertemplate{navigation symbols}{}
\setbeamertemplate{itemize items}[circle]

% Frame title bar
\setbeamertemplate{frametitle}{%
  \nointerlineskip
  \begin{beamercolorbox}[wd=\paperwidth,ht=3.1ex,dp=1.2ex,leftskip=1.2em]{frametitle}
    \usebeamerfont{frametitle}\insertframetitle
  \end{beamercolorbox}%
  \nointerlineskip
  \hbox{\color{gold}\rule{\paperwidth}{1.2pt}}%
}

% Footer
\setbeamertemplate{footline}{%
  \leavevmode%
  \hbox{%
    \begin{beamercolorbox}[wd=.72\paperwidth,ht=2.4ex,dp=1ex,left,leftskip=1em]{footline}
      \color{softgray}\scriptsize Universidad de Pamplona $\cdot$ Faculty of Basic Sciences $\cdot$ Computational Physics II
    \end{beamercolorbox}%
    \begin{beamercolorbox}[wd=.28\paperwidth,ht=2.4ex,dp=1ex,right,rightskip=1em]{footline}
      \color{softgray}\scriptsize \insertframenumber\,/\,\inserttotalframenumber
    \end{beamercolorbox}}%
  \vskip0pt%
}

% ---------------------------------------------------------------------------
% tcolorbox styles (institutional look)
% ---------------------------------------------------------------------------
\newtcolorbox{keybox}[1]{colback=winelight,colframe=wine,coltitle=white,
  fonttitle=\bfseries\small,title={#1},boxrule=0.6pt,arc=2pt,
  left=4pt,right=4pt,top=2pt,bottom=2pt,before skip=3pt,after skip=3pt}
\newtcolorbox{warnbox}[1]{colback=gold!10,colframe=gold,coltitle=white,
  fonttitle=\bfseries\small,title={#1},boxrule=0.6pt,arc=2pt,
  left=4pt,right=4pt,top=2pt,bottom=2pt,before skip=3pt,after skip=3pt}

% Math shortcuts
\newcommand{\ii}{\mathrm{i}}
\newcommand{\ee}{\mathrm{e}}
\newcommand{\dd}{\,\mathrm{d}}

\title{Fourier Series and the Discrete Fourier Transform}
\subtitle{Computational Physics II \,--\, Unit 4: Fourier Transform I}
\author{Cristian Vargas}
\institute{Universidad de Pamplona \\ Faculty of Basic Sciences \\ Physics Program}
\date{\today}

\begin{document}

% ===========================================================================
\begin{frame}[plain]
  \begin{tikzpicture}[remember picture,overlay]
    \fill[wine] (current page.north west) rectangle ([yshift=-0.46\paperheight]current page.north east);
    \fill[gold] ([yshift=-0.46\paperheight]current page.north west) rectangle ([yshift=-0.475\paperheight]current page.north east);
  \end{tikzpicture}
  \vspace{0.6cm}
  \begin{center}
    {\color{white}\LARGE\bfseries Fourier Series and\\[2pt] the Discrete Fourier Transform}\\[10pt]
    {\color{white}\large Computational Physics II \,--\, Unit 4: Fourier Transform I}
  \end{center}
  \vspace{1.3cm}
  \begin{center}
    {\Large Cristian Vargas}\\[4pt]
    Universidad de Pamplona \,$\cdot$\, Faculty of Basic Sciences\\[2pt]
    \today
  \end{center}
\end{frame}

% ===========================================================================
\begin{frame}{Contents}
  \footnotesize
  \begin{columns}[T]
    \column{0.5\textwidth}
    \textbf{Part A. Fourier series}
    \begin{enumerate}
      \item Physical motivation
      \item Orthogonality of sines and cosines
      \item Computing the coefficients
      \item Symmetries and simplifications
      \item Example 1: the square wave
      \item Example 2: the triangle wave
      \item Convergence and the Gibbs phenomenon
      \item Smoothness and coefficient decay
      \item Arbitrary period
      \item Complex form
      \item Parseval's theorem
      \item Application: driven oscillator
      \item Application: heat equation
    \end{enumerate}
    \column{0.5\textwidth}
    \textbf{Part B. Discrete Fourier Transform}
    \begin{enumerate}\setcounter{enumi}{13}
      \item From the series to the DFT
      \item Definition and discrete orthogonality
      \item The inverse transform
      \item A worked example ($N=4$)
      \item Frequencies, Nyquist and aliasing
      \item Physical amplitude and phase
      \item Cost and the FFT
      \item Spectral leakage and windows
      \item Zero padding
      \item Python code and results
      \item Analysis checklist and summary
    \end{enumerate}
  \end{columns}
\end{frame}

% ===========================================================================
\begin{frame}{1. Motivation: why Fourier in physics?}
  \begin{itemize}
    \item Many physical systems are \textbf{linear}: waves, RLC circuits, oscillators, heat diffusion, the free Schr\"odinger equation.
    \item In a linear system \textbf{superposition} holds: if we know the response to $\cos(\omega t)$ and $\sin(\omega t)$, we know the response to any combination of them.
    \item The exponentials $\ee^{\ii\omega t}$ are \textbf{eigenfunctions} of the derivative operator:
    \[
      \frac{\mathrm{d}}{\mathrm{d}t}\,\ee^{\ii\omega t}=\ii\omega\,\ee^{\ii\omega t}.
    \]
    So a linear differential equation becomes an \emph{algebraic} equation for each frequency.
    \item \textbf{Central idea:} decompose an arbitrary periodic function into a sum of harmonic oscillations, each with its own amplitude and phase.
  \end{itemize}
  \begin{keybox}{Goal of this presentation}
    Understand where every formula comes from (not just use it), justify every simplification, and take everything to a numerical code with a physical application.
  \end{keybox}
\end{frame}

% ===========================================================================
\begin{frame}{2. Orthogonality of sines and cosines}
  \small
  We work with functions of period $2\pi$ on $[-\pi,\pi]$ (we generalize later).\\[3pt]
  \textbf{Result} (for integers $m,n\ge 1$):
  \[
    \int_{-\pi}^{\pi}\!\cos mx\cos nx\dd x=\pi\,\delta_{mn},\quad
    \int_{-\pi}^{\pi}\!\sin mx\sin nx\dd x=\pi\,\delta_{mn},\quad
    \int_{-\pi}^{\pi}\!\sin mx\cos nx\dd x=0 .
  \]
  \textbf{Proof.} We use the product-to-sum identity:
  \[
    \cos mx\cos nx=\tfrac12\big[\cos (m-n)x+\cos (m+n)x\big].
  \]
  For an integer $k\neq 0$:
  $\displaystyle\int_{-\pi}^{\pi}\cos kx\dd x=\frac{2\sin k\pi}{k}=0$; for $k=0$ it equals $2\pi$.
  \begin{itemize}
    \item If $m\neq n$: both $m-n\neq0$ and $m+n\neq 0$, so both integrals vanish $\Rightarrow$ result $0$.
    \item If $m=n$: $\tfrac12[2\pi+0]=\pi$.
  \end{itemize}
  For $\sin\cdot\cos$ the integrand is \textbf{odd} (odd times even), and the integral of an odd function over a symmetric interval is zero.
  \begin{keybox}{Interpretation}
    $\{1,\cos nx,\sin nx\}$ is an orthogonal basis of the space of periodic functions, with inner product $\langle f,g\rangle=\int_{-\pi}^{\pi}f g\dd x$.
  \end{keybox}
\end{frame}

% ===========================================================================
\begin{frame}{3. Computing the Fourier coefficients}
  \footnotesize
  Assume $f$ can be written as
  \[
    f(x)=\frac{a_0}{2}+\sum_{n=1}^{\infty}\big[a_n\cos nx+b_n\sin nx\big].
  \]
  \textbf{To isolate $a_m$:} multiply by $\cos mx$ and integrate over $[-\pi,\pi]$:
  \[
    \int_{-\pi}^{\pi}\! f\cos mx\dd x=\frac{a_0}{2}\underbrace{\int\cos mx\dd x}_{0\ (m\ge1)}
    +\sum_n\Big[a_n\underbrace{\int\cos nx\cos mx\dd x}_{\pi\delta_{nm}}+b_n\underbrace{\int\sin nx\cos mx\dd x}_{0}\Big]=\pi a_m .
  \]
  \begin{columns}[T]
    \column{0.58\textwidth}
    \[
      \boxed{a_n=\frac1\pi\!\int_{-\pi}^{\pi}\! f\cos nx\dd x,\qquad
      b_n=\frac1\pi\!\int_{-\pi}^{\pi}\! f\sin nx\dd x}
    \]
    \column{0.4\textwidth}
    For $m=0$: $\int f\dd x=\tfrac{a_0}{2}\,2\pi\Rightarrow a_0=\tfrac1\pi\!\int f\dd x$.
  \end{columns}
  \begin{keybox}{Why $a_0/2$ and not $a_0$?}
    For $m=0$ the integral of $\cos^2$ is $2\pi$, not $\pi$. Writing the constant term as $a_0/2$ makes \emph{the same formula} $a_n=\frac1\pi\int f\cos nx$ valid for every $n\ge0$. Also, $a_0/2$ is the \textbf{mean value} of $f$.
  \end{keybox}
  \begin{warnbox}{Assumption made}
    Integrating the series term by term: valid if the series converges uniformly (or, in general, in the $L^2$ sense).
  \end{warnbox}
\end{frame}

% ===========================================================================
\begin{frame}{4. Symmetries: when coefficients vanish}
  \small
  Basic rule: $\displaystyle\int_{-\pi}^{\pi}g(x)\dd x=0$ if $g$ is odd, and $=2\int_0^{\pi}g(x)\dd x$ if $g$ is even.
  \begin{center}
  \begin{tabular}{@{}lll@{}}
    \toprule
    If $f$ is\ldots & $f\cos nx$ is\ldots & $f\sin nx$ is\ldots\\
    \midrule
    \textbf{even} & even $\Rightarrow$ $a_n=\frac2\pi\!\int_0^\pi f\cos nx$ & odd $\Rightarrow b_n=0$\\
    \textbf{odd}  & odd $\Rightarrow a_n=0$ & even $\Rightarrow b_n=\frac2\pi\!\int_0^\pi f\sin nx$\\
    \bottomrule
  \end{tabular}
  \end{center}
  \begin{keybox}{Why this simplification is legitimate}
    It is not an approximation, it is an exact equality. It follows from (even)$\times$(even)=even, (odd)$\times$(odd)=even, (even)$\times$(odd)=odd, plus the change of variable $x\to-x$ on the negative half of the interval: $\int_{-\pi}^{0}g(x)\dd x=\int_{0}^{\pi}g(-x)\dd x$.
  \end{keybox}
  \textbf{Practical benefits:}
  \begin{itemize}
    \item The integration interval is halved (less work).
    \item You know \emph{in advance} which coefficients must be zero: an excellent \textbf{check} of numerical results.
    \item An extra symmetry ($f(x+\pi)=-f(x)$, ``half-wave'') makes \textbf{only odd harmonics survive}.
  \end{itemize}
\end{frame}

% ===========================================================================
\begin{frame}{5. Example 1: the square wave (step by step)}
  \small
  \[
    f(x)=\begin{cases}-1,&-\pi<x<0\\ +1,&0<x<\pi\end{cases}\qquad f(x+2\pi)=f(x).
  \]
  \begin{enumerate}
    \item $f$ is \textbf{odd} $\Rightarrow a_n=0$ for all $n\ge0$ (including $a_0$: zero mean value).
    \item Coefficients $b_n$ (by symmetry we integrate only on $[0,\pi]$, where $f=1$):
    \[
      b_n=\frac2\pi\int_0^\pi\sin nx\dd x=\frac2\pi\Big[-\frac{\cos nx}{n}\Big]_0^\pi
      =\frac{2\,(1-\cos n\pi)}{n\pi}=\frac{2\,[1-(-1)^n]}{n\pi}.
    \]
    \item Since $(-1)^n=-1$ for odd $n$ and $+1$ for even $n$:
    \[
      b_n=\begin{cases}\dfrac{4}{n\pi},& n\text{ odd}\\[4pt] 0,& n\text{ even}\end{cases}
      \quad\Longrightarrow\quad
      f(x)=\frac4\pi\sum_{k=1}^{\infty}\frac{\sin\big((2k-1)x\big)}{2k-1}.
    \]
  \end{enumerate}
  \begin{keybox}{Check: Leibniz series}
    At $x=\pi/2$ the series gives $1=\frac4\pi\big(1-\frac13+\frac15-\cdots\big)$, i.e. $\frac\pi4=1-\frac13+\frac15-\cdots$ (a known result). Even harmonics are missing because $f(x+\pi)=-f(x)$.
  \end{keybox}
\end{frame}

% ===========================================================================
\begin{frame}{6. Example 2: the triangle wave $f(x)=|x|$}
  \small
  On $[-\pi,\pi]$, $f(x)=|x|$, extended with period $2\pi$. It is \textbf{even} $\Rightarrow b_n=0$.
  \[
    a_0=\frac2\pi\int_0^\pi x\dd x=\pi,\qquad
    a_n=\frac2\pi\int_0^\pi x\cos nx\dd x=\frac2\pi\Big[\frac{x\sin nx}{n}+\frac{\cos nx}{n^2}\Big]_0^\pi=\frac{2\,[(-1)^n-1]}{\pi n^2}.
  \]
  (Integration by parts with $u=x$, $\dd v=\cos nx\dd x$.) Hence
  \[
    a_n=\begin{cases}-\dfrac{4}{\pi n^2},& n\text{ odd}\\[4pt]0,& n\text{ even}\end{cases}
    \qquad\Longrightarrow\qquad
    |x|=\frac\pi2-\frac4\pi\sum_{k=1}^{\infty}\frac{\cos\big((2k-1)x\big)}{(2k-1)^2}.
  \]
  \begin{columns}[T]
    \column{0.52\textwidth}
    \begin{keybox}{Two independent checks}
      \begin{itemize}\footnotesize
        \item $x=0$: $0=\frac\pi2-\frac4\pi\sum\frac1{(2k-1)^2}=\frac\pi2-\frac4\pi\cdot\frac{\pi^2}{8}=0$ \checkmark
        \item Parseval: $\frac{\pi^2}{4}+\frac12\sum a_n^2=\frac{\pi^2}{4}+\frac{\pi^2}{12}=\frac{\pi^2}{3}=\frac1{2\pi}\!\int x^2$ \checkmark
      \end{itemize}
    \end{keybox}
    \column{0.45\textwidth}
    \textbf{Key difference with the square wave:} here $f$ is \emph{continuous}; only its derivative jumps. The coefficients decay as $1/n^2$ (instead of $1/n$) and there is \textbf{no Gibbs overshoot}.
  \end{columns}
\end{frame}

% ===========================================================================
\begin{frame}{7. Convergence and the Gibbs phenomenon}
  \small
  \textbf{Dirichlet conditions} (sufficient): within one period $f$ is bounded, has finitely many maxima/minima and finitely many finite discontinuities. Then the series converges to $f(x)$ at continuity points and to the \textbf{average of the one-sided limits} at a jump (here $0$ at $x=0$).
  \begin{warnbox}{Gibbs: the overshoot does not go away}
    Near a jump of size $\Delta$, partial sums overshoot the limit by $\approx 9\%\,\Delta$ \emph{no matter how many terms are used}; it only becomes narrower.
  \end{warnbox}
  \textbf{Derivation (partial sum with $K$ odd harmonics).}
  \[
    S_K'(x)=\frac4\pi\sum_{k=1}^{K}\cos(2k-1)x=\frac2\pi\,\frac{\sin 2Kx}{\sin x}
    \quad\Rightarrow\quad\text{first maximum at }x_K=\frac{\pi}{2K}.
  \]
  \[
    S_K(x_K)=\frac2\pi\int_0^{x_K}\frac{\sin 2Ku}{\sin u}\dd u
    \;\xrightarrow{u=v/2K,\ \sin u\approx u}\;
    \frac2\pi\int_0^{\pi}\frac{\sin v}{v}\dd v=\frac2\pi\,\mathrm{Si}(\pi)\approx1.1790 .
  \]
  \footnotesize\textbf{Why $\sin u\approx u$ is justified:} in the integration range $u\le\pi/2K$, which is tiny for large $K$, so the relative error is $O(u^2)\to0$. That is why the limit is exact as $K\to\infty$ and independent of $K$.
\end{frame}

% ===========================================================================
\begin{frame}{Figure: partial sums of the square wave}
  \centering
  \includegraphics[height=0.74\textheight]{fig1_partial_sums.png}\\[2pt]
  \small With $N=49$ the approximation is good away from the jump, but the overshoot near $x=0$ is still $\approx 0.179$.
\end{frame}

% ===========================================================================
\begin{frame}{8. Smoothness and coefficient decay}
  \footnotesize
  \textbf{Claim.} If $f$ and its first $p-1$ derivatives are continuous and $f^{(p)}$ is piecewise continuous, then $c_n=O(1/n^{p+1})$.
  \textbf{Derivation} (integration by parts, $T$-periodic $f$, $\omega_0=2\pi/T$):
  \[
    c_n=\frac1T\int_0^T\! f\,\ee^{-\ii n\omega_0t}\dd t
    =\underbrace{\frac1T\Big[\frac{f\,\ee^{-\ii n\omega_0 t}}{-\ii n\omega_0}\Big]_0^T}_{=0\text{ if }f\text{ continuous \& periodic}}+\frac{1}{\ii n\omega_0}\underbrace{\frac1T\int_0^T\! f'\,\ee^{-\ii n\omega_0 t}\dd t}_{c_n(f')}.
  \]
  \begin{columns}[T]
    \column{0.47\textwidth}
    \begin{itemize}
      \item Each continuous derivative contributes one factor $1/n$.
      \item Square wave: $p=0\Rightarrow 1/n$.
      \item Triangle wave: $p=1\Rightarrow1/n^2$.
      \item Analytic functions: faster than any power (exponential decay).
    \end{itemize}
    \begin{keybox}{Practical meaning}
      \footnotesize Smoother $\Rightarrow$ fewer harmonics needed. With 5 harmonics ($N=9$) the triangle wave error is $0.064$ everywhere; for the square wave it is $0.182$ even away from the jumps.
    \end{keybox}
    \column{0.52\textwidth}
    \includegraphics[width=\textwidth]{fig4_decay_rates.png}
  \end{columns}
\end{frame}

% ===========================================================================
\begin{frame}{9. Arbitrary period $T$}
  If $f(t+T)=f(t)$, define $\omega_0=2\pi/T$ and change variable $x=\omega_0 t$ (so $x\in[-\pi,\pi]$ when $t\in[-T/2,T/2]$).
  \[
    f(t)=\frac{a_0}{2}+\sum_{n=1}^{\infty}\big[a_n\cos(n\omega_0 t)+b_n\sin(n\omega_0 t)\big]
  \]
  \[
    a_n=\frac2T\int_{-T/2}^{T/2}\! f(t)\cos(n\omega_0 t)\dd t,\qquad
    b_n=\frac2T\int_{-T/2}^{T/2}\! f(t)\sin(n\omega_0 t)\dd t .
  \]
  \begin{keybox}{Where does the factor $2/T$ come from?}
    With $x=\omega_0t$, $\dd x=\omega_0\dd t$, so $\frac1\pi\dd x=\frac{\omega_0}{\pi}\dd t=\frac2T\dd t$.
  \end{keybox}
  \textbf{Physical meaning:} harmonic $n$ oscillates at frequency $f_n=n/T$ (fundamental $f_1=1/T$). A long period means harmonics very close in frequency: $\Delta f=1/T$. This relation will return as the \textbf{frequency resolution} of the DFT.
\end{frame}

% ===========================================================================
\begin{frame}{10. Complex form of the series}
  \small
  Using Euler: $\cos\theta=\tfrac12(\ee^{\ii\theta}+\ee^{-\ii\theta})$, $\ \sin\theta=\tfrac1{2\ii}(\ee^{\ii\theta}-\ee^{-\ii\theta})$. Then
  \[
    a_n\cos n\omega_0t+b_n\sin n\omega_0t=\underbrace{\frac{a_n-\ii b_n}{2}}_{c_n}\ee^{\ii n\omega_0 t}+\underbrace{\frac{a_n+\ii b_n}{2}}_{c_{-n}}\ee^{-\ii n\omega_0 t}.
  \]
  \[
    \boxed{f(t)=\sum_{n=-\infty}^{\infty}c_n\,\ee^{\ii n\omega_0 t},\qquad
    c_n=\frac1T\int_{0}^{T}f(t)\,\ee^{-\ii n\omega_0 t}\dd t}
  \]
  \begin{itemize}
    \item \textbf{Complex orthogonality:} $\frac1T\int_0^T\ee^{\ii (m-n)\omega_0 t}\dd t=\delta_{mn}$. $c_n$ is isolated as before, multiplying by $\ee^{-\ii n\omega_0t}$.
    \item Relations: $c_0=a_0/2$, $\ a_n=2\,\mathrm{Re}\,c_n$, $\ b_n=-2\,\mathrm{Im}\,c_n$ (for $n\ge1$).
    \item \textbf{If $f$ is real}, $c_{-n}=c_n^{*}$: the negative spectrum carries no new information.
  \end{itemize}
  \begin{keybox}{Why use the complex form?}
    It is more compact (one formula), amplitude and phase sit in a single complex number $c_n=|c_n|\ee^{\ii\varphi_n}$, and it generalizes directly to the Fourier transform and the DFT.
  \end{keybox}
\end{frame}

% ===========================================================================
\begin{frame}{11. Parseval's theorem}
  \small
  Compute the ``mean power'' of a real $f$: $\frac1T\int_0^T|f|^2\dd t$. Substituting the complex series for $f$ and for $f^*$:
  \[
    \frac1T\int_0^T\! f f^{*}\dd t=\sum_{m,n}c_m c_n^{*}\underbrace{\frac1T\int_0^T\!\ee^{\ii(m-n)\omega_0t}\dd t}_{\delta_{mn}}
    =\sum_{n=-\infty}^{\infty}|c_n|^2=\frac{a_0^2}{4}+\frac12\sum_{n=1}^{\infty}\big(a_n^2+b_n^2\big).
  \]
  \begin{keybox}{Physical interpretation}
    The (average) energy of the signal is the sum of the energies of each harmonic: there are no ``cross terms'' because harmonics are orthogonal. The factor $\frac12$ is the average of $\cos^2$ or $\sin^2$ over a period.
  \end{keybox}
  \textbf{Check with the square wave} ($|f|^2=1$, so the left side equals 1):
  \[
    \frac12\sum_{n\ \text{odd}}\frac{16}{n^2\pi^2}=\frac{8}{\pi^2}\sum_{n\ \text{odd}}\frac1{n^2}=\frac8{\pi^2}\cdot\frac{\pi^2}{8}=1\ \checkmark
  \]
  With $N=9$ harmonics we already recover $96\%$ of the power; with $N=999$, $99.96\%$ (result of the code).
\end{frame}

% ===========================================================================
\begin{frame}{12. Application: damped driven oscillator}
  \small
  \[
    m\ddot x+\gamma\dot x+kx=F(t),\qquad F(t+T)=F(t),\quad \omega_0^2=k/m,\ \ \beta=\gamma/m,\ \ \omega=\tfrac{2\pi}{T}.
  \]
  Write $F(t)=\sum_n F_n\,\ee^{\ii n\omega t}$ and look for $x(t)=\sum_n x_n\,\ee^{\ii n\omega t}$ (steady state). Since $\frac{\mathrm d}{\mathrm dt}\ee^{\ii n\omega t}=\ii n\omega\,\ee^{\ii n\omega t}$:
  \[
    m\big(\omega_0^2-n^2\omega^2+\ii\beta n\omega\big)\,x_n=F_n
    \quad\Longrightarrow\quad
    x_n=\frac{F_n/m}{\omega_0^2-n^2\omega^2+\ii\beta n\omega}.
  \]
  \begin{itemize}
    \item \textbf{Justification:} the equation is linear and the exponentials are linearly independent (orthogonal), so the coefficient of each $\ee^{\ii n\omega t}$ is matched separately.
    \item Each harmonic is amplified by the \textbf{response function} $\chi(n\omega)=1/(\omega_0^2-n^2\omega^2+\ii\beta n\omega)$: resonance when $n\omega\approx\omega_0$.
    \item For a square-wave force only odd harmonics appear ($F_n\propto1/n$). If $3\omega\approx\omega_0$, the \textbf{third harmonic} excites the resonance even though the fundamental frequency is far away.
  \end{itemize}
\end{frame}

% ===========================================================================
\begin{frame}{13. Application: heat equation on a rod}
  \small
  Rod of length $L$ with ends held at $0$: $\ \partial_t u=\alpha\,\partial_x^2u$, $\ u(0,t)=u(L,t)=0$, $\ u(x,0)=f(x)$.
  \textbf{Separation of variables} $u=X(x)T(t)$: $\ X''=-\lambda X$ with $X(0)=X(L)=0$ forces
  \[
    X_n(x)=\sin\frac{n\pi x}{L},\quad\lambda_n=\Big(\frac{n\pi}{L}\Big)^2,\qquad T_n(t)=\ee^{-\alpha\lambda_n t}.
  \]
  Superposition and the initial condition give a \textbf{Fourier sine series} (period $2L$, odd extension of $f$):
  \[
    u(x,t)=\sum_{n=1}^{\infty}b_n\,\sin\frac{n\pi x}{L}\;\ee^{-\alpha(n\pi/L)^2t},\qquad
    b_n=\frac2L\int_0^L f(x)\sin\frac{n\pi x}{L}\dd x .
  \]
  \begin{columns}[T]
    \column{0.55\textwidth}
    \begin{keybox}{Why a sine series?}
      \footnotesize The boundary conditions $u=0$ at both ends select the sines (they vanish at $0$ and $L$); the cosines are discarded. Orthogonality of the sines on $[0,L]$ gives $b_n$.
    \end{keybox}
    \column{0.43\textwidth}
    \textbf{Physics:} mode $n$ decays with rate $\propto n^2$. High harmonics (sharp features, like Gibbs ripples) vanish first: \textbf{diffusion smooths}.
  \end{columns}
\end{frame}

% ===========================================================================
\begin{frame}{14. From the series to the DFT: discretize the integral}
  \small
  In practice we do not know $f(t)$ continuously, only \textbf{samples} (measurements or numerical results):
  \[
    x_m=f(t_m),\qquad t_m=m\,\Delta t=\frac{mT}{N},\qquad m=0,1,\dots,N-1,\qquad f_s=\frac1{\Delta t}=\frac NT .
  \]
  Approximate $c_k=\frac1T\int_0^T f(t)\,\ee^{-\ii k\omega_0t}\dd t$ with the rectangle rule:
  \[
    c_k\approx\frac1T\sum_{m=0}^{N-1}x_m\,\ee^{-\ii k\frac{2\pi}{T}\frac{mT}{N}}\,\Delta t
    =\frac1N\sum_{m=0}^{N-1}x_m\,\ee^{-2\pi\ii\,km/N}\equiv\frac{X_k}{N}.
  \]
  \begin{keybox}{Why is the simplest rule good enough?}
    The integrand is \textbf{periodic}: its values at $t=0$ and $t=T$ coincide. For periodic, smooth integrands the rectangle rule equals the trapezoid rule, and the endpoint error terms of the Euler--Maclaurin formula cancel between the two ends. The error decays faster than any power of $1/N$ (spectral convergence).
  \end{keybox}
  This finite sum is, by definition, the \textbf{DFT}: $\ X_k=\sum_{m=0}^{N-1}x_m\,\ee^{-2\pi\ii km/N}$.
\end{frame}

% ===========================================================================
\begin{frame}{15. Definition and discrete orthogonality}
  \footnotesize
  Let $W_N=\ee^{-2\pi\ii/N}$ (primitive $N$-th root of unity: $W_N^N=1$).
  \[
    \boxed{X_k=\sum_{m=0}^{N-1}x_m\,W_N^{km},\qquad k=0,1,\dots,N-1}
  \]
  \textbf{Lemma (discrete orthogonality).} For integer $p$:
  \[
    \sum_{m=0}^{N-1}\ee^{2\pi\ii\,pm/N}=\begin{cases}N,&p\equiv0\pmod N\\0,&\text{otherwise.}\end{cases}
  \]
  \textbf{Proof.} If $p\equiv0$, each term is 1 and there are $N$ of them. Otherwise, with $r=\ee^{2\pi\ii p/N}\neq1$ it is a geometric series:
  \[
    \sum_{m=0}^{N-1}r^m=\frac{1-r^N}{1-r}=\frac{1-\ee^{2\pi\ii p}}{1-r}=0,
    \quad\text{because }\ee^{2\pi\ii p}=1\text{ for integer }p.
  \]
  \begin{keybox}{Meaning}
    The $N$ vectors $\{\ee^{2\pi\ii km/N}\}_{k=0}^{N-1}$ are orthogonal in $\mathbb C^N$: the DFT is a \textbf{change of basis} (a unitary rotation up to the factor $\sqrt N$), just like the Fourier series is in function space.
  \end{keybox}
\end{frame}

% ===========================================================================
\begin{frame}{16. The inverse transform}
  \small
  \[
    \boxed{x_m=\frac1N\sum_{k=0}^{N-1}X_k\,\ee^{+2\pi\ii\,km/N}}
  \]
  \textbf{Proof.} Substitute the definition of $X_k$ into the right-hand side (dummy index $n$ in the inner sum):
  \[
    \frac1N\sum_{k=0}^{N-1}\Big[\sum_{n=0}^{N-1}x_n\,\ee^{-2\pi\ii kn/N}\Big]\ee^{2\pi\ii km/N}
    =\sum_{n=0}^{N-1}x_n\;\underbrace{\frac1N\sum_{k=0}^{N-1}\ee^{2\pi\ii k(m-n)/N}}_{\delta_{mn}\ \text{(lemma with }p=m-n)}=x_m .
  \]
  Since $0\le m,n\le N-1$, $m-n\equiv0\ (\mathrm{mod}\ N)$ only when $m=n$.
  \begin{warnbox}{Normalization conventions}
    There are different choices for the $1/N$ factor: (i) in the inverse (NumPy, SciPy); (ii) $1/\sqrt N$ in both (unitary transform); (iii) in the forward transform. All are equivalent as long as the \emph{same} one is used in both directions. Here: NumPy's, $X_k=\sum x_m \ee^{-2\pi\ii km/N}$.
  \end{warnbox}
  \textbf{Matrix form:} $\mathbf X=F\,\mathbf x$ with $F_{km}=W_N^{km}$ and $F^{-1}=\tfrac1N F^{\dagger}$.
\end{frame}

% ===========================================================================
\begin{frame}{17. A worked example: $N=4$ by hand}
  \small
  Take $\mathbf x=(1,\,2,\,0,\,-1)$. Here $W_4=\ee^{-\ii\pi/2}=-\ii$, so $W_4^0=1,\ W_4^1=-\ii,\ W_4^2=-1,\ W_4^3=\ii$ and
  \[
    F=\begin{pmatrix}1&1&1&1\\1&-\ii&-1&\ii\\1&-1&1&-1\\1&\ii&-1&-\ii\end{pmatrix}\!,\qquad
    \mathbf X=F\mathbf x=\begin{pmatrix}1+2+0-1\\ 1-2\ii+0-\ii\\ 1-2+0+1\\ 1+2\ii+0+\ii\end{pmatrix}=\begin{pmatrix}2\\1-3\ii\\0\\1+3\ii\end{pmatrix}.
  \]
  \begin{columns}[T]
    \column{0.5\textwidth}
    \begin{keybox}{Checks}
      \footnotesize
      \begin{itemize}
        \item \textbf{Real signal:} $X_3=X_1^{*}$ \checkmark\ (and $X_0,X_2$ real).
        \item \textbf{Mean:} $X_0/N=0.5=\bar x$ \checkmark
        \item \textbf{Discrete Parseval:} $\sum|x_m|^2=6$ and $\frac1N\sum|X_k|^2=\frac{4+10+0+10}{4}=6$ \checkmark
      \end{itemize}
    \end{keybox}
    \column{0.47\textwidth}
    \textbf{Inverse:} e.g. $x_1=\tfrac14\big(2+(1-3\ii)\ii+0+(1+3\ii)(-\ii)\big)=\tfrac14\,(2+3+\ii+3-\ii)=2$ \checkmark\\[3pt]
    \footnotesize The code (Part 8) reproduces this and gives zero difference with \texttt{np.fft.fft}.
  \end{columns}
\end{frame}

% ===========================================================================
\begin{frame}{18. Frequencies, Nyquist and aliasing}
  \small
  Since $k\omega_0=2\pi k/T$ and $T=N/f_s$, index $k$ corresponds to the frequency
  \[
    f_k=\frac{k}{T}=k\,\frac{f_s}{N},\qquad\Delta f=\frac{f_s}{N}=\frac1T\quad\text{(resolution).}
  \]
  \textbf{Periodicity in $k$:} $\ee^{-2\pi\ii(k+N)m/N}=\ee^{-2\pi\ii km/N}\,\ee^{-2\pi\ii m}=\ee^{-2\pi\ii km/N}$ $\Rightarrow$ $X_{k+N}=X_k$. Only $N$ independent values.
  \textbf{Real signal:} $X_{N-k}=\sum x_m\ee^{-2\pi\ii(N-k)m/N}=\sum x_m\ee^{+2\pi\ii km/N}=X_k^{*}$.
  \begin{columns}[T]
    \column{0.52\textwidth}
    \begin{keybox}{Nyquist theorem}
      \footnotesize With sampling frequency $f_s$ only frequencies $|f|<f_s/2$ (index $k<N/2$) can be represented. A component at $f>f_s/2$ appears \emph{disguised} (alias) at $f-nf_s$.
    \end{keybox}
    \column{0.46\textwidth}
    \textbf{Example:} with $f_s=100$ Hz, a $93$ Hz signal looks like $|93-100|=7$ Hz.\\[2pt]
    \textbf{Why:} $\cos(2\pi\cdot93\,t_m)=\cos(2\pi\cdot7\,t_m)$ at the instants $t_m=m/100$.
  \end{columns}
  \begin{itemize}
    \item \textbf{Rule of thumb:} choose $f_s>2f_{\max}$ (ideally $5$--$10\,f_{\max}$) and filter before sampling.
    \item In NumPy: \texttt{np.fft.rfft} (only $k=0..N/2$, for real signals) and \texttt{np.fft.rfftfreq(N, d=1/fs)}.
  \end{itemize}
\end{frame}

% ===========================================================================
\begin{frame}{19. Physical amplitude and phase from the DFT}
  \small
  Let $x_m=A\cos(2\pi f_0t_m+\varphi)$ with $f_0=k_0f_s/N$ (\emph{exact multiple} of $\Delta f$, see spectral leakage). Then
  \[
    x_m=\frac A2\Big[\ee^{\ii\varphi}\ee^{2\pi\ii k_0m/N}+\ee^{-\ii\varphi}\ee^{-2\pi\ii k_0m/N}\Big].
  \]
  Applying the DFT and the orthogonality lemma (each exponential only ``activates'' one $k$):
  \[
    X_{k_0}=\frac{NA}{2}\,\ee^{\ii\varphi},\qquad X_{N-k_0}=\frac{NA}{2}\,\ee^{-\ii\varphi},\qquad\text{all other }X_k=0 .
  \]
  \begin{keybox}{Formulas to read a spectrum}
    \[
      A=\frac{2}{N}\,|X_{k_0}|,\qquad\varphi=\arg X_{k_0}\qquad(k_0\neq0,\ N/2)
    \]
    The factor $2$ accounts for the energy being split between $+f_0$ and $-f_0$. For $k=0$ (DC component) the mean value is $X_0/N$ (no factor 2).
  \end{keybox}
  This is exactly what the code does in Part 5: it recovers $A_1\approx1.005$, $A_2\approx0.486$ and phases $0.284$ and $1.106$ rad (true values: $1$, $0.5$, $0.3$, $1.1$), with added noise of amplitude $0.3$.
\end{frame}

% ===========================================================================
\begin{frame}{Figure: normal modes of two coupled oscillators}
  \centering
  \includegraphics[height=0.78\textheight]{fig2_normal_modes.png}\\[2pt]
  \small Noisy signal in time (top) and its spectrum (bottom): sharp peaks at 3 Hz and 7 Hz.
\end{frame}

% ===========================================================================
\begin{frame}{20. Cost and the FFT}
  \small
  Evaluating $X_k=\sum_mx_mW_N^{km}$ for every $k$ costs $N^2$ complex multiplications: $O(N^2)$.
  \textbf{Cooley--Tukey idea} ($N$ even): split $m$ into even ($m=2r$) and odd ($m=2r+1$):
  \[
    X_k=\underbrace{\sum_{r=0}^{N/2-1}x_{2r}\,W_{N/2}^{kr}}_{E_k}+W_N^{k}\underbrace{\sum_{r=0}^{N/2-1}x_{2r+1}\,W_{N/2}^{kr}}_{O_k},
    \qquad\text{since }W_N^{2kr}=W_{N/2}^{kr}.
  \]
  $E_k$ and $O_k$ are size-$N/2$ DFTs, periodic in $k$ with period $N/2$. Also $W_N^{N/2}=-1$, so
  \[
    X_k=E_k+W_N^kO_k,\qquad X_{k+N/2}=E_k-W_N^kO_k\qquad(k=0,\dots,N/2-1).
  \]
  \begin{keybox}{Result}
    Recurrence $T(N)=2T(N/2)+O(N)\Rightarrow T(N)=O(N\log_2N)$. For $N=4096$: $\sim5\times10^4$ operations vs $\sim1.7\times10^7$. In the code (Part 4): \textbf{0.59 s} (hand-written DFT) vs \textbf{0.14 ms} (FFT), numerical difference $\sim10^{-10}$ (round-off).
  \end{keybox}
  The simplest version requires $N=2^p$; libraries handle any $N$ (slower if $N$ has large prime factors).
\end{frame}

% ===========================================================================
\begin{frame}{21. Spectral leakage and windows}
  \small
  The DFT implicitly assumes the sampled signal \textbf{repeats periodically} with period $N\Delta t$.
  \begin{itemize}
    \item If the window contains a \emph{whole} number of periods, the repetition is continuous and the spectrum is a clean delta.
    \item If not ($f_0$ is not a multiple of $f_s/N$), the repetition has an \textbf{artificial jump} at the edges, and energy ``leaks'' into many neighboring frequencies (like Gibbs: a jump $\Rightarrow$ coefficients decaying as $1/k$).
  \end{itemize}
  Mathematically: observing $N$ samples equals multiplying by a rectangular window, and a product in time is a \textbf{convolution} in frequency with the kernel $\sin(N\pi\Delta f)/\sin(\pi\Delta f)$, whose side lobes decay slowly.
  \begin{keybox}{Hann window: $w_m=\tfrac12\big[1-\cos(2\pi m/N)\big]$}
    It takes the signal smoothly to zero at both ends, removing the jump. Cost: the main peak widens ($\approx2\times$). Normalize by $\sum w_m$ instead of $N$.
    In the code (Part 6): leakage far from the peak drops from $1.7\times10^{-2}$ to $1.7\times10^{-5}$ (a factor of a thousand).
  \end{keybox}
  \footnotesize\textbf{Rule of thumb:} if possible, measure a whole number of periods; if not, apply a window.
\end{frame}

% ===========================================================================
\begin{frame}{Figure: leakage with and without a window}
  \centering
  \includegraphics[height=0.72\textheight]{fig3_spectral_leakage.png}\\[2pt]
  \small $f_0=3.37$ Hz is not a multiple of $\Delta f=0.25$ Hz. The Hann window drastically reduces the side lobes.
\end{frame}

% ===========================================================================
\begin{frame}{22. Zero padding: finer grid, same resolution}
  \small
  Observing a signal for a time $T$ gives a \emph{continuous} spectrum (the DTFT) whose features have width $\sim1/T$. The DFT only \textbf{samples} it at $f_k=kf_s/N$.
  \begin{itemize}
    \item Appending zeros ($N\to N_{\rm pad}$) samples the same spectrum on a finer grid $f_s/N_{\rm pad}$: \textbf{interpolation}, no new information.
    \item The \textbf{true resolution} is set by $T=N/f_s$: two tones closer than $\sim1/T$ stay merged, whatever the padding.
    \item Useful for: locating a peak more precisely, smooth plots, and making $N_{\rm pad}=2^p$ (fast FFT).
  \end{itemize}
  \begin{columns}[T]
    \column{0.42\textwidth}
    \begin{keybox}{Example in the code (Part 9)}
      \footnotesize Tones at $10$ and $10.6$ Hz, $T=1$ s ($1/T=1$ Hz). With $N_{\rm pad}=100$ or $800$ they appear as a single peak: padding only draws it smoothly.
    \end{keybox}
    \begin{warnbox}{To resolve them}
      \footnotesize Measure longer. Numerically (Hann window), $T=4$ s already shows two peaks at $10.0$ and $10.6$ Hz, while $T=1$ s cannot.
    \end{warnbox}
    \column{0.56\textwidth}
    \includegraphics[width=\textwidth]{fig5_zero_padding.png}
  \end{columns}
\end{frame}

% ===========================================================================
\begin{frame}{23. The Python code: structure}
  \small
  File: \texttt{fourier\_physics.py} (each part is an independent, commented function).
  \begin{center}
  \footnotesize
  \begin{tabular}{@{}p{0.06\textwidth}p{0.46\textwidth}p{0.40\textwidth}@{}}
    \toprule
    \textbf{Part} & \textbf{What it does} & \textbf{Theory concept}\\
    \midrule
    1 & Analytic $a_n,b_n$ vs numerical integration (trapezoid) & Orthogonality, odd symmetry\\
    2 & Partial sums $S_N(t)$ and overshoot & Convergence, Gibbs\\
    3 & $\frac12\sum b_n^2\to1$ & Parseval\\
    4 & DFT as matrix $F\mathbf x$ vs \texttt{np.fft.fft} & Definition, $O(N^2)$ vs $O(N\log N)$\\
    5 & Two normal modes + noise; read $f,A,\varphi$ & Frequencies, physical amplitude\\
    6 & Non-integer $f_0$, no window vs Hann & Spectral leakage\\
    7 & Square vs triangle wave decay, Parseval & Smoothness vs decay\\
    8 & $N=4$ by hand vs code & Definition, symmetries\\
    9 & Zero padding of two close tones & Resolution vs grid\\
    \bottomrule
  \end{tabular}
  \end{center}
  \begin{itemize}\footnotesize
    \item \textbf{Run:} \texttt{python fourier\_physics.py} (needs NumPy and Matplotlib). It saves 5 PNG figures.
    \item \textbf{Version note:} in NumPy $\ge2.0$ \texttt{np.trapz} was renamed \texttt{np.trapezoid}; the code supports both.
    \item In Part 1, $a_n\sim-10^{-5}$ instead of exactly $0$: an endpoint artifact of the trapezoid rule (discretization slightly breaks periodicity at $t=\pm\pi$), not a flaw of the theory.
  \end{itemize}
\end{frame}

% ===========================================================================
\begin{frame}[fragile]{24. Key code excerpts}
  \footnotesize
  \textbf{DFT as a matrix product} (definition, $O(N^2)$):
\begin{verbatim}
def dft_manual(x):
    N = len(x)
    m = np.arange(N)
    k = m.reshape((N, 1))
    F = np.exp(-2j * np.pi * k * m / N)   # F_km = W_N^(km)
    return F @ x                          # X = F x
\end{verbatim}
  \textbf{Spectrum of a real signal with physical amplitudes:}
\begin{verbatim}
X     = np.fft.rfft(x)                   # k = 0 ... N/2
freqs = np.fft.rfftfreq(N, d=1/fs)       # f_k = k * fs / N
amp   = 2.0 * np.abs(X) / N              # A = 2|X_k|/N
amp[0] /= 2.0                            # DC term has no factor 2
phase = np.angle(X)                      # phi = arg X_k
\end{verbatim}
  \textbf{Hann window (against leakage):}
\begin{verbatim}
hann = 0.5 * (1 - np.cos(2*np.pi*np.arange(N)/N))
amp_w = 2 * np.abs(np.fft.rfft(x * hann)) / np.sum(hann)
\end{verbatim}
\end{frame}

% ===========================================================================
\begin{frame}{25. Numerical results obtained}
  \footnotesize
  \begin{columns}[T]
    \column{0.49\textwidth}
    \textbf{Part 1 -- coefficients} ($b_n$)
    \begin{center}
    \begin{tabular}{@{}rcc@{}}
      \toprule $n$ & analytic & numeric\\\midrule
      1 & 1.27323954 & 1.27323954\\
      3 & 0.42441318 & 0.42441318\\
      5 & 0.25464791 & 0.25464791\\
      7 & 0.18189136 & 0.18189136\\
      9 & 0.14147106 & 0.14147106\\
      \bottomrule
    \end{tabular}
    \end{center}
    \textbf{Part 2 -- Gibbs overshoot}
    \begin{center}
    \begin{tabular}{@{}rc@{}}
      \toprule $N$ & $\max S_N-1$\\\midrule
      9 & 0.182328\\ 49 & 0.179113\\ 199 & 0.178988\\ 999 & 0.178980\\
      \bottomrule
    \end{tabular}\\[2pt]
    Theory: $\frac2\pi\mathrm{Si}(\pi)-1=0.178980$
    \end{center}
    \column{0.49\textwidth}
    \textbf{Part 3 -- Parseval (square wave)}
    \begin{center}
    \begin{tabular}{@{}rc@{}}
      \toprule $N$ & $\frac12\sum b_n^2$\\\midrule
      9 & 0.95960\\ 99 & 0.99595\\ 999 & 0.99959\\ 9999 & 0.99996\\
      \bottomrule
    \end{tabular}
    \end{center}
    \textbf{Part 5 -- normal modes}
    \begin{center}
    \begin{tabular}{@{}cccc@{}}
      \toprule mode & $f$ [Hz] & $A$ & $\varphi$ [rad]\\\midrule
      1 (true: 1, 0.3) & 3.00 & 1.005 & 0.284\\
      2 (true: 0.5, 1.1) & 7.00 & 0.486 & 1.106\\
      \bottomrule
    \end{tabular}
    \end{center}
    \textbf{Part 7 -- triangle Parseval:} $3.28986813$ vs $\pi^2/3=3.28986813$.\\
    \textbf{Part 8:} hand-computed $N=4$ DFT equals \texttt{np.fft.fft} exactly.
  \end{columns}
\end{frame}

% ===========================================================================
\begin{frame}{26. Analysis checklist for a real signal}
  \small
  \begin{enumerate}
    \item \textbf{Choose $f_s$} at least $2.5$--$10$ times the highest frequency of interest (Nyquist); low-pass filter if needed.
    \item \textbf{Choose the duration $T$} from the required resolution: $\Delta f\approx1/T$ (and, with a window, about $2\times$ coarser).
    \item \textbf{Remove the mean (and trend)}, otherwise a huge $k=0$ peak leaks into nearby bins.
    \item \textbf{Apply a window} (Hann) unless an integer number of periods is guaranteed.
    \item Compute \texttt{rfft} and \texttt{rfftfreq}; convert with $A=2|X_k|/N$ (divide by $\sum w_m$ if windowed).
    \item \textbf{Zero-pad} (optional) for smoother plots or precise peak location; remember it does not improve resolution.
    \item \textbf{Interpret:} peaks $\leftrightarrow$ normal modes / driving harmonics; check that nothing sits near $f_s/2$ (possible aliasing).
  \end{enumerate}
  \begin{keybox}{Typical physics uses}
    Normal-mode frequencies of oscillators, vibration spectra, diffraction and optics (Fourier optics), solving linear PDEs, filtering experimental noise, and spectral analysis of simulation output.
  \end{keybox}
\end{frame}

% ===========================================================================
\begin{frame}{27. Summary and common mistakes}
  \footnotesize
  \begin{columns}[T]
    \column{0.52\textwidth}
    \begin{keybox}{Key ideas}
      \begin{itemize}
        \item Fourier series: expansion in an orthogonal basis; coefficients come from projecting.
        \item Symmetries of $f$ remove coefficients \emph{exactly}.
        \item Discontinuity $\Rightarrow$ $c_n\sim1/n$ and Gibbs; smoother $f\Rightarrow$ faster decay.
        \item DFT = rectangle rule on the complex series: $c_k\approx X_k/N$.
        \item Resolution $\Delta f=f_s/N=1/T$; maximum frequency $f_s/2$.
        \item FFT: same result as the DFT, in $O(N\log N)$.
      \end{itemize}
    \end{keybox}
    \column{0.46\textwidth}
    \begin{warnbox}{Common mistakes}
      \begin{itemize}
        \item Forgetting the $2/N$ factor for physical amplitude.
        \item Reading $k$ as a frequency (it is $f_k=kf_s/N$).
        \item Sampling below $2f_{\max}$ (aliasing).
        \item Ignoring leakage: window with non-integer periods.
        \item Thinking zero padding improves resolution.
        \item Mixing normalization conventions between libraries.
        \item Reading $k>N/2$ as high frequencies (they are \emph{negative} frequencies).
      \end{itemize}
    \end{warnbox}
  \end{columns}
  \vspace{2pt}
  \footnotesize\textbf{Suggested references:} M. Newman, \emph{Computational Physics}, ch. 7 (Fourier transforms); M. Boas, \emph{Mathematical Methods in the Physical Sciences}, ch. 7; J. W. Cooley and J. W. Tukey (1965), \emph{Math. Comput.} 19, 297.
\end{frame}

\end{document}
